% !Mode:: "TeX:UTF-8"

\chapter{主体内容}
% 本章示范图、表、公式、算法、代码等常用格式的使用方法

\section{本章引言}

每章开头可用简短的引言段落说明本章主要研究内容，以及与前后章节之间的逻辑关系。

\section{图的使用}

\subsection{单幅图}

图序及图名居中置于图的下方，采用五号楷体字。图、表、公式一律采用阿拉伯数字分章编号，例如图\ref{fig:example}所示。

\begin{figure}[htbp]
    \centering
    \includegraphics[width=0.6\textwidth]{example-image} % 请替换为实际图片
    \caption{示例图片标题}
    \label{fig:example}
\end{figure}

\subsection{双语图标题}

如需添加中英文双语图标题，可使用 \verb|\bilingualfigure| 命令，如图\ref{fig:bilingual}所示。

\begin{figure}[htbp]
    \centering
    \includegraphics[width=0.6\textwidth]{example-image}
    \bilingualfigure{双语标题示例图}{Example figure with bilingual caption}
    \label{fig:bilingual}
\end{figure}

\subsection{子图排列}

包含多个子图时，可使用 \texttt{subfigure} 环境，如图\ref{fig:subfig}所示。

\begin{figure}[htbp]
    \centering
    \begin{subfigure}[b]{0.45\textwidth}
        \centering
        \includegraphics[width=\linewidth]{example-image-a}
        \caption{子图 (a)}
    \end{subfigure}
    \hspace{10pt}
    \begin{subfigure}[b]{0.45\textwidth}
        \centering
        \includegraphics[width=\linewidth]{example-image-b}
        \caption{子图 (b)}
    \end{subfigure}
    \caption{子图排列示例}
    \label{fig:subfig}
\end{figure}

\section{表的使用}

表格采用国际通用三线表，表序及表名置于表的上方，采用五号楷体字。如表\ref{tab:example}所示。

\begin{table}[htbp]
    \centering
    \caption{示例三线表}
    \label{tab:example}
    \begin{tabular}{cccc}
        \toprule
        方法 & 准确率 (\%) & 召回率 (\%) & F1 分数 \\
        \midrule
        方法 A & 92.3 & 89.1 & 90.7 \\
        方法 B & 90.5 & 91.2 & 90.8 \\
        本文方法 & \textbf{94.6} & \textbf{93.8} & \textbf{94.2} \\
        \bottomrule
    \end{tabular}
\end{table}

\section{公式的使用}

公式编号用括号括起写在右边行末，格式为"章号-公式序号"，如式（\ref{eq:example}）所示。

\begin{equation}
    \mathcal{L} = -\frac{1}{N}\sum_{i=1}^{N} \left[ y_i \log \hat{y}_i + (1-y_i)\log(1-\hat{y}_i) \right]
    \label{eq:example}
\end{equation}

行内公式示例：设 $x \in \mathbb{R}^d$ 为输入特征，$W \in \mathbb{R}^{d \times k}$ 为权重矩阵，则输出为 $y = Wx + b$。

多行对齐公式示例：

\begin{align}
    f(x) &= ax^2 + bx + c \label{eq:poly} \\
    f'(x) &= 2ax + b \label{eq:derive}
\end{align}

\section{定理与证明}

\begin{definition}[定义名称]
    这里给出定义的内容。
\end{definition}

\begin{theorem}[定理名称]
    这里给出定理的内容和命题。
\end{theorem}

\begin{proof}
    这里给出定理的证明过程。证明过程应逻辑严谨，步骤清晰。
\end{proof}

\section{算法描述}

使用 \texttt{algorithm2e} 宏包描述算法，如算法\ref{alg:example}所示。

\begin{algorithm}[H]
    \caption{示例算法}
    \label{alg:example}
    \KwIn{输入数据 $\mathcal{D}$，学习率 $\eta$，迭代次数 $T$}
    \KwOut{模型参数 $\theta$}
    初始化参数 $\theta_0$\;
    \For{$t = 1, 2, \ldots, T$}{
        计算梯度 $g_t = \nabla_\theta \mathcal{L}(\theta_{t-1}; \mathcal{D})$\;
        更新参数 $\theta_t = \theta_{t-1} - \eta \cdot g_t$\;
        \If{收敛条件满足}{
            \textbf{break}\;
        }
    }
    \Return{$\theta_T$}
\end{algorithm}

\section{代码示例}

程序代码使用等宽字体，并配有行号，如代码\ref{lst:example}所示。

\begin{lstlisting}[language=Python, caption={Python 代码示例}, label={lst:example}]
import numpy as np

def softmax(x):
    """计算 softmax 函数"""
    e_x = np.exp(x - np.max(x))
    return e_x / e_x.sum(axis=0)

# 示例调用
logits = np.array([2.0, 1.0, 0.5])
probs = softmax(logits)
print(probs)
\end{lstlisting}

\section{本章小结}

每章末尾应有小结，对本章研究工作进行概括性总结，并说明所取得的主要成果。同时可指出本章工作的局限性，为下文进一步研究做铺垫。
